% Options for packages loaded elsewhere
\PassOptionsToPackage{unicode}{hyperref}
\PassOptionsToPackage{hyphens}{url}
\documentclass[
  11pt,
  a4paper,
]{article}
\usepackage{xcolor}
\usepackage[margin=18mm]{geometry}
\usepackage{amsmath,amssymb}
\setcounter{secnumdepth}{-\maxdimen} % remove section numbering
\usepackage{iftex}
\ifPDFTeX
  \usepackage[T1]{fontenc}
  \usepackage[utf8]{inputenc}
  \usepackage{textcomp} % provide euro and other symbols
\else % if luatex or xetex
  \usepackage{unicode-math} % this also loads fontspec
  \defaultfontfeatures{Scale=MatchLowercase}
  \defaultfontfeatures[\rmfamily]{Ligatures=TeX,Scale=1}
\fi
\ifPDFTeX\else
  % xetex/luatex font selection
    \setmainfont[]{Noto Serif}
    \setsansfont[]{Noto Sans}
    \setmonofont[]{Noto Sans Mono}
  \setmathfont[]{Noto Sans Math}
\fi
% Use upquote if available, for straight quotes in verbatim environments
\IfFileExists{upquote.sty}{\usepackage{upquote}}{}
\IfFileExists{microtype.sty}{% use microtype if available
  \usepackage[]{microtype}
  \UseMicrotypeSet[protrusion]{basicmath} % disable protrusion for tt fonts
}{}
\makeatletter
\@ifundefined{KOMAClassName}{% if non-KOMA class
  \IfFileExists{parskip.sty}{%
    \usepackage{parskip}
  }{% else
    \setlength{\parindent}{0pt}
    \setlength{\parskip}{6pt plus 2pt minus 1pt}}
}{% if KOMA class
  \KOMAoptions{parskip=half}}
\makeatother
\ifLuaTeX
\usepackage[bidi=basic]{babel}
\else
\usepackage[bidi=default]{babel}
\fi
\babelprovide[main,import]{russian}
\ifPDFTeX
\else
\babelfont{rm}[]{Noto Serif}
\fi
% get rid of language-specific shorthands (see #6817):
\let\LanguageShortHands\languageshorthands
\def\languageshorthands#1{}
\setlength{\emergencystretch}{3em} % prevent overfull lines
\providecommand{\tightlist}{%
  \setlength{\itemsep}{0pt}\setlength{\parskip}{0pt}}
\usepackage{bookmark}
\IfFileExists{xurl.sty}{\usepackage{xurl}}{} % add URL line breaks if available
\urlstyle{same}
\hypersetup{
  pdftitle={Билеты по дискретным случайным процессам},
  pdflang={ru-RU},
  hidelinks,
  pdfcreator={LaTeX via pandoc}}

\title{Билеты по дискретным случайным процессам}
\author{}
\date{}

\begin{document}
\maketitle

\section{Билет 20. Конечномерные распределения как обобщенные функции.
Дельта-функции Дирака и преобразования
Фурье}\label{ux431ux438ux43bux435ux442-20.-ux43aux43eux43dux435ux447ux43dux43eux43cux435ux440ux43dux44bux435-ux440ux430ux441ux43fux440ux435ux434ux435ux43bux435ux43dux438ux44f-ux43aux430ux43a-ux43eux431ux43eux431ux449ux435ux43dux43dux44bux435-ux444ux443ux43dux43aux446ux438ux438.-ux434ux435ux43bux44cux442ux430-ux444ux443ux43dux43aux446ux438ux438-ux434ux438ux440ux430ux43aux430-ux438-ux43fux440ux435ux43eux431ux440ux430ux437ux43eux432ux430ux43dux438ux44f-ux444ux443ux440ux44cux435}

Источники: Ширяев 1, §12; лекции ДСП.

\subsection{1. Короткий
ответ}\label{ux43aux43eux440ux43eux442ux43aux438ux439-ux43eux442ux432ux435ux442}

Конечномерное распределение процесса можно рассматривать не только как
меру на \(\mathbb R^n\), но и как линейный функционал на функциях. Как
мера оно действует на ограниченные борелевские или непрерывные функции;
если говорить строго на языке обобщенных функций, его ограничивают на
тестовые функции из \(C_c^\infty\) или на функции Шварца. Если

\[
\mu=P_{t_1,\ldots,t_n}
\]

\begin{itemize}
\tightlist
\item
  распределение вектора
\end{itemize}

\[
(X_{t_1},\ldots,X_{t_n}),
\]

то каждой ограниченной непрерывной функции \(f\in C_b(\mathbb R^n)\) оно
сопоставляет число

\[
\langle \mu,f\rangle
=
\int_{\mathbb R^n} f(x)\,\mu(dx)
=
\mathbb E f(X_{t_1},\ldots,X_{t_n}).
\]

В этом смысле мера действует как обобщенная функция.

Дельта-мера Дирака в точке \(a\):

\[
\delta_a(B)=\mathbf 1_B(a),
\]

а как функционал:

\[
\langle\delta_a,f\rangle=f(a).
\]

Преобразование Фурье вероятностной меры:

\[
\widehat\mu(u)
=
\int_{\mathbb R^n} e^{i\langle u,x\rangle}\,\mu(dx).
\]

Для распределения случайного вектора это его характеристическая функция:

\[
\varphi_X(u)=\mathbb E e^{i\langle u,X\rangle}.
\]

\subsection{2. Полный
ответ}\label{ux43fux43eux43bux43dux44bux439-ux43eux442ux432ux435ux442}

\subsubsection{2.1. Вероятностная мера и конечномерные
распределения}\label{ux432ux435ux440ux43eux44fux442ux43dux43eux441ux442ux43dux430ux44f-ux43cux435ux440ux430-ux438-ux43aux43eux43dux435ux447ux43dux43eux43cux435ux440ux43dux44bux435-ux440ux430ux441ux43fux440ux435ux434ux435ux43bux435ux43dux438ux44f}

Этот билет связывает несколько уже изученных тем: распределения
случайных векторов из Билета 06, характеристические функции из Билета
08, моменты через характеристические функции из Билета 13 и меру Дирака
из Билета 01.

Обычно распределение случайного вектора определяется как вероятностная
мера. Пусть

\[
X=(X_1,\ldots,X_n)
\]

\begin{itemize}
\tightlist
\item
  случайный вектор. Его распределение
\end{itemize}

\[
\mu=P_X
\]

на \(\mathbb R^n\) задается формулой

\[
\mu(B)=P\{X\in B\},
\qquad
B\in\mathcal B(\mathbb R^n).
\]

Для случайного процесса конечномерное распределение - это распределение
вектора

\[
(X_{t_1},\ldots,X_{t_n}).
\]

Его обычно обозначают

\[
P_{t_1,\ldots,t_n}
\]

или

\[
\mu_{t_1,\ldots,t_n}.
\]

\subsubsection{2.2. Распределение как функционал на пространстве
тестовых
функций}\label{ux440ux430ux441ux43fux440ux435ux434ux435ux43bux435ux43dux438ux435-ux43aux430ux43a-ux444ux443ux43dux43aux446ux438ux43eux43dux430ux43b-ux43dux430-ux43fux440ux43eux441ux442ux440ux430ux43dux441ux442ux432ux435-ux442ux435ux441ux442ux43eux432ux44bux445-ux444ux443ux43dux43aux446ux438ux439}

Но меру можно задавать не только значениями на множествах, а и ее
действием на функции. Если \(f\) - ограниченная борелевская или
ограниченная непрерывная функция, то можно взять интеграл

\[
\int_{\mathbb R^n} f(x)\,\mu(dx).
\]

Для вероятностных мер особенно удобен класс ограниченных непрерывных
тестовых функций

\[
C_b(\mathbb R^n).
\]

Ограниченность гарантирует, что интеграл конечен:

\[
\left|\int f\,d\mu\right|
\le
\|f\|_\infty\mu(\mathbb R^n)
=
\|f\|_\infty.
\]

Непрерывность важна для слабой сходимости распределений.

Таким образом, распределение \(\mu\) задает функционал

\[
f\mapsto \langle \mu,f\rangle,
\]

где

\[
\langle \mu,f\rangle
=
\int f(x)\,\mu(dx).
\]

\subsubsection{2.3. Свойства функционала и интерпретация как обобщенной
функции}\label{ux441ux432ux43eux439ux441ux442ux432ux430-ux444ux443ux43dux43aux446ux438ux43eux43dux430ux43bux430-ux438-ux438ux43dux442ux435ux440ux43fux440ux435ux442ux430ux446ux438ux44f-ux43aux430ux43a-ux43eux431ux43eux431ux449ux435ux43dux43dux43eux439-ux444ux443ux43dux43aux446ux438ux438}

Этот функционал линейный:

\[
\langle \mu,af+bg\rangle
=
a\langle\mu,f\rangle+b\langle\mu,g\rangle.
\]

Он положительный:

\[
f\ge0
\quad\Rightarrow\quad
\langle\mu,f\rangle\ge0.
\]

Он нормирован:

\[
\langle\mu,1\rangle=1.
\]

Именно поэтому можно говорить, что распределение рассматривается как
обобщенная функция: оно действует на тестовую функцию и возвращает
число. Здесь слово ``обобщенная'' не означает, что это обязательно
обычная функция плотности. Например, атомарное распределение не имеет
плотности относительно меры Лебега, но как функционал на тестовых
функциях оно определено без проблем.

Для случайного вектора \(X\) действие распределения на тестовую функцию
можно записать через математическое ожидание:

\[
\langle P_X,f\rangle
=
\int_{\mathbb R^n}f(x)\,P_X(dx)
=
\mathbb E f(X).
\]

Для конечномерного распределения процесса:

\[
\langle P_{t_1,\ldots,t_n},f\rangle
=
\mathbb E f(X_{t_1},\ldots,X_{t_n}).
\]

Это особенно удобно: вместо работы с вероятностями всех борелевских
множеств можно проверять интегралы от тестовых функций.

\subsubsection{2.4. Мера Дирака и способы представления
распределений}\label{ux43cux435ux440ux430-ux434ux438ux440ux430ux43aux430-ux438-ux441ux43fux43eux441ux43eux431ux44b-ux43fux440ux435ux434ux441ux442ux430ux432ux43bux435ux43dux438ux44f-ux440ux430ux441ux43fux440ux435ux434ux435ux43bux435ux43dux438ux439}

Теперь дельта-функция Дирака. Строго в теории меры это не обычная
функция, а мера, сосредоточенная в одной точке. Для \(a\in\mathbb R^n\):

\[
\delta_a(B)=
\begin{cases}
1, & a\in B,\\
0, & a\notin B.
\end{cases}
\]

Эквивалентно:

\[
\delta_a(B)=\mathbf 1_B(a).
\]

Как функционал дельта действует подстановкой:

\[
\langle\delta_a,f\rangle
=
\int f(x)\,\delta_a(dx)
=
f(a).
\]

Физическая запись

\[
\int f(x)\delta(x-a)\,dx=f(a)
\]

понимается именно в этом смысле. Строго \(\delta(x-a)\) не является
обычной функцией плотности по Лебегу.

Если случайная величина постоянна:

\[
X=a
\quad\text{п.н.},
\]

то ее распределение равно

\[
P_X=\delta_a.
\]

Если распределение дискретное:

\[
P\{X=x_k\}=p_k,
\qquad
\sum_k p_k=1,
\]

то его можно записать как сумму дельта-мер:

\[
P_X=\sum_k p_k\delta_{x_k}.
\]

Действие на тестовую функцию:

\[
\langle P_X,f\rangle
=
\sum_k p_k f(x_k).
\]

Если распределение абсолютно непрерывно с плотностью \(p(x)\), то

\[
\mu(dx)=p(x)\,dx
\]

и

\[
\langle\mu,f\rangle
=
\int_{\mathbb R^n}f(x)p(x)\,dx.
\]

Так в одной записи объединяются дискретные, непрерывные и смешанные
распределения.

\subsubsection{2.5. Слабая сходимость
распределений}\label{ux441ux43bux430ux431ux430ux44f-ux441ux445ux43eux434ux438ux43cux43eux441ux442ux44c-ux440ux430ux441ux43fux440ux435ux434ux435ux43bux435ux43dux438ux439}

Следующая важная идея - слабая сходимость. Последовательность
вероятностных мер \(\mu_k\) слабо сходится к \(\mu\), если для всех
ограниченных непрерывных функций \(f\):

\[
\int f\,d\mu_k\to\int f\,d\mu.
\]

Пишут:

\[
\mu_k\Rightarrow\mu.
\]

Для случайных величин это означает сходимость по распределению:

\[
X_k\Rightarrow X.
\]

Иными словами, слабая сходимость - это сходимость распределений как
функционалов на тестовых функциях.

\subsubsection{2.6. Преобразование Фурье меры и характеристические
функции}\label{ux43fux440ux435ux43eux431ux440ux430ux437ux43eux432ux430ux43dux438ux435-ux444ux443ux440ux44cux435-ux43cux435ux440ux44b-ux438-ux445ux430ux440ux430ux43aux442ux435ux440ux438ux441ux442ux438ux447ux435ux441ux43aux438ux435-ux444ux443ux43dux43aux446ux438ux438}

Теперь преобразование Фурье меры. Для конечной меры \(\mu\) на
\(\mathbb R^n\) ее преобразование Фурье определяется формулой

\[
\widehat\mu(u)
=
\int_{\mathbb R^n}e^{i\langle u,x\rangle}\,\mu(dx),
\qquad
u\in\mathbb R^n.
\]

Так как

\[
|e^{i\langle u,x\rangle}|=1,
\]

интеграл всегда существует для вероятностной меры.

Если \(\mu=P_X\) - распределение случайного вектора \(X\), то

\[
\widehat\mu(u)
=
\mathbb E e^{i\langle u,X\rangle}
=
\varphi_X(u).
\]

То есть характеристическая функция - это просто преобразование Фурье
распределения.

Для дельта-мер:

\[
\widehat{\delta_a}(u)
=
\int e^{i\langle u,x\rangle}\delta_a(dx)
=
e^{i\langle u,a\rangle}.
\]

Для дискретного распределения:

\[
\mu=\sum_k p_k\delta_{x_k}
\]

получаем

\[
\widehat\mu(u)=\sum_k p_k e^{i\langle u,x_k\rangle}.
\]

Для распределения с плотностью:

\[
\widehat\mu(u)=\int_{\mathbb R^n}e^{i\langle u,x\rangle}p(x)\,dx.
\]

\subsubsection{2.7. Свертка и теорема
единственности}\label{ux441ux432ux435ux440ux442ux43aux430-ux438-ux442ux435ux43eux440ux435ux43cux430-ux435ux434ux438ux43dux441ux442ux432ux435ux43dux43dux43eux441ux442ux438}

Преобразование Фурье удобно тем, что свертка мер превращается в
произведение. Если \(X\) и \(Y\) независимы, то распределение суммы
\(X+Y\) равно свертке распределений:

\[
P_{X+Y}=P_X*P_Y,
\]

а характеристические функции перемножаются:

\[
\varphi_{X+Y}(u)=\varphi_X(u)\varphi_Y(u).
\]

Это использовалось в Билете 08 и Билете 19.

Еще одно ключевое свойство: характеристическая функция однозначно задает
распределение. Если

\[
\widehat\mu(u)=\widehat\nu(u)
\]

для всех \(u\), то

\[
\mu=\nu.
\]

Поэтому преобразование Фурье - не просто удобный числовой образ, а
полноценный способ кодировать распределение.

\subsection{3. Основные
определения}\label{ux43eux441ux43dux43eux432ux43dux44bux435-ux43eux43fux440ux435ux434ux435ux43bux435ux43dux438ux44f}

\subsubsection{Конечномерное распределение
процесса}\label{ux43aux43eux43dux435ux447ux43dux43eux43cux435ux440ux43dux43eux435-ux440ux430ux441ux43fux440ux435ux434ux435ux43bux435ux43dux438ux435-ux43fux440ux43eux446ux435ux441ux441ux430}

Для процесса \((X_t)\) и набора \(t_1,\ldots,t_n\) это распределение
случайного вектора

\[
(X_{t_1},\ldots,X_{t_n})
\]

на \(\mathbb R^n\):

\[
P_{t_1,\ldots,t_n}(B)
=
P\{(X_{t_1},\ldots,X_{t_n})\in B\}.
\]

\subsubsection{Тестовая
функция}\label{ux442ux435ux441ux442ux43eux432ux430ux44f-ux444ux443ux43dux43aux446ux438ux44f}

В этом билете удобно брать

\[
f\in C_b(\mathbb R^n),
\]

то есть \(f\) непрерывна и ограничена.

\subsubsection{Мера как
функционал}\label{ux43cux435ux440ux430-ux43aux430ux43a-ux444ux443ux43dux43aux446ux438ux43eux43dux430ux43b}

Вероятностная мера \(\mu\) действует на тестовую функцию так:

\[
\langle\mu,f\rangle=\int f\,d\mu.
\]

\subsubsection{Дельта-мера
Дирака}\label{ux434ux435ux43bux44cux442ux430-ux43cux435ux440ux430-ux434ux438ux440ux430ux43aux430}

Мера в точке \(a\):

\[
\delta_a(B)=\mathbf 1_B(a).
\]

Как функционал:

\[
\langle\delta_a,f\rangle=f(a).
\]

\subsubsection{Дискретное распределение через
дельта-меры}\label{ux434ux438ux441ux43aux440ux435ux442ux43dux43eux435-ux440ux430ux441ux43fux440ux435ux434ux435ux43bux435ux43dux438ux435-ux447ux435ux440ux435ux437-ux434ux435ux43bux44cux442ux430-ux43cux435ux440ux44b}

\[
\mu=\sum_k p_k\delta_{x_k}.
\]

\subsubsection{Абсолютно непрерывное
распределение}\label{ux430ux431ux441ux43eux43bux44eux442ux43dux43e-ux43dux435ux43fux440ux435ux440ux44bux432ux43dux43eux435-ux440ux430ux441ux43fux440ux435ux434ux435ux43bux435ux43dux438ux435}

Если есть плотность \(p\), то

\[
\mu(dx)=p(x)\,dx.
\]

\subsubsection{Слабая сходимость
мер}\label{ux441ux43bux430ux431ux430ux44f-ux441ux445ux43eux434ux438ux43cux43eux441ux442ux44c-ux43cux435ux440}

\[
\mu_n\Rightarrow\mu
\quad\Longleftrightarrow\quad
\int f\,d\mu_n\to\int f\,d\mu
\]

для всех \(f\in C_b\).

\subsubsection{Преобразование Фурье
меры}\label{ux43fux440ux435ux43eux431ux440ux430ux437ux43eux432ux430ux43dux438ux435-ux444ux443ux440ux44cux435-ux43cux435ux440ux44b}

\[
\widehat\mu(u)=\int e^{i\langle u,x\rangle}\,\mu(dx).
\]

\subsubsection{Характеристическая
функция}\label{ux445ux430ux440ux430ux43aux442ux435ux440ux438ux441ux442ux438ux447ux435ux441ux43aux430ux44f-ux444ux443ux43dux43aux446ux438ux44f}

Для \(X\sim\mu\):

\[
\varphi_X(u)=\mathbb E e^{i\langle u,X\rangle}=\widehat\mu(u).
\]

\subsection{4. Основные теоремы и
свойства}\label{ux43eux441ux43dux43eux432ux43dux44bux435-ux442ux435ux43eux440ux435ux43cux44b-ux438-ux441ux432ux43eux439ux441ux442ux432ux430}

\begin{itemize}
\tightlist
\item
  Вероятностная мера задает положительный нормированный линейный
  функционал:
\end{itemize}

\[
f\mapsto \int f\,d\mu.
\]

\begin{itemize}
\tightlist
\item
  Для конечномерного распределения:
\end{itemize}

\[
\int f\,dP_{t_1,\ldots,t_n}
=
\mathbb E f(X_{t_1},\ldots,X_{t_n}).
\]

\begin{itemize}
\tightlist
\item
  Дельта-мера действует как подстановка:
\end{itemize}

\[
\int f\,d\delta_a=f(a).
\]

\begin{itemize}
\tightlist
\item
  Дискретное распределение записывается как сумма атомов:
\end{itemize}

\[
\mu=\sum_kp_k\delta_{x_k}.
\]

\begin{itemize}
\tightlist
\item
  Характеристическая функция - преобразование Фурье распределения:
\end{itemize}

\[
\varphi_X(u)=\widehat{P_X}(u).
\]

\begin{itemize}
\tightlist
\item
  Преобразование Фурье дельта-меры:
\end{itemize}

\[
\widehat{\delta_a}(u)=e^{i\langle u,a\rangle}.
\]

\begin{itemize}
\tightlist
\item
  Свертка переходит в произведение:
\end{itemize}

\[
\widehat{\mu*\nu}(u)=\widehat\mu(u)\widehat\nu(u).
\]

\begin{itemize}
\tightlist
\item
  Характеристическая функция однозначно задает распределение.
\item
  Слабая сходимость распределений проверяется на ограниченных
  непрерывных тестовых функциях.
\end{itemize}

\subsection{5. Примеры}\label{ux43fux440ux438ux43cux435ux440ux44b}

\subsubsection{Постоянная случайная
величина}\label{ux43fux43eux441ux442ux43eux44fux43dux43dux430ux44f-ux441ux43bux443ux447ux430ux439ux43dux430ux44f-ux432ux435ux43bux438ux447ux438ux43dux430}

Если

\[
X=a
\quad\text{п.н.},
\]

то

\[
P_X=\delta_a.
\]

Для тестовой функции:

\[
\mathbb E f(X)=f(a).
\]

Характеристическая функция:

\[
\varphi_X(t)=e^{ita}.
\]

\subsubsection{Дискретное
распределение}\label{ux434ux438ux441ux43aux440ux435ux442ux43dux43eux435-ux440ux430ux441ux43fux440ux435ux434ux435ux43bux435ux43dux438ux435}

Пусть

\[
P\{X=x_k\}=p_k.
\]

Тогда

\[
P_X=\sum_kp_k\delta_{x_k}.
\]

Действие на тестовую функцию:

\[
\langle P_X,f\rangle=\sum_kp_kf(x_k).
\]

Характеристическая функция:

\[
\varphi_X(t)=\sum_kp_ke^{itx_k}.
\]

\subsubsection{Бернуллиевская
величина}\label{ux431ux435ux440ux43dux443ux43bux43bux438ux435ux432ux441ux43aux430ux44f-ux432ux435ux43bux438ux447ux438ux43dux430}

Если

\[
P\{X=1\}=p,
\qquad
P\{X=0\}=1-p,
\]

то

\[
P_X=(1-p)\delta_0+p\delta_1.
\]

Характеристическая функция:

\[
\varphi_X(t)=1-p+pe^{it}.
\]

\subsubsection{Распределение с
плотностью}\label{ux440ux430ux441ux43fux440ux435ux434ux435ux43bux435ux43dux438ux435-ux441-ux43fux43bux43eux442ux43dux43eux441ux442ux44cux44e}

Если

\[
\mu(dx)=p(x)\,dx,
\]

то

\[
\langle\mu,f\rangle=\int f(x)p(x)\,dx.
\]

Преобразование Фурье:

\[
\widehat\mu(t)=\int e^{itx}p(x)\,dx.
\]

\subsubsection{Конечномерное распределение
процесса}\label{ux43aux43eux43dux435ux447ux43dux43eux43cux435ux440ux43dux43eux435-ux440ux430ux441ux43fux440ux435ux434ux435ux43bux435ux43dux438ux435-ux43fux440ux43eux446ux435ux441ux441ux430-1}

Для процесса \((X_n)\) двумерное распределение пары \((X_0,X_1)\)
действует на функцию \(f(x,y)\) так:

\[
\langle P_{0,1},f\rangle
=
\mathbb E f(X_0,X_1).
\]

Если процесс детерминированный и

\[
X_0=a,\qquad X_1=b
\]

почти наверное, то

\[
P_{0,1}=\delta_{(a,b)}.
\]

\subsection{6. Схемы доказательств /
обоснований}\label{ux441ux445ux435ux43cux44b-ux434ux43eux43aux430ux437ux430ux442ux435ux43bux44cux441ux442ux432-ux43eux431ux43eux441ux43dux43eux432ux430ux43dux438ux439}

\subsubsection{Почему мера задает линейный положительный
функционал}\label{ux43fux43eux447ux435ux43cux443-ux43cux435ux440ux430-ux437ux430ux434ux430ux435ux442-ux43bux438ux43dux435ux439ux43dux44bux439-ux43fux43eux43bux43eux436ux438ux442ux435ux43bux44cux43dux44bux439-ux444ux443ux43dux43aux446ux438ux43eux43dux430ux43b}

Для \(f,g\in C_b\) и чисел \(a,b\):

\[
\int(af+bg)\,d\mu
=
a\int f\,d\mu+b\int g\,d\mu.
\]

Если \(f\ge0\), то

\[
\int f\,d\mu\ge0.
\]

Если \(\mu\) - вероятность, то

\[
\int 1\,d\mu=\mu(\mathbb R^n)=1.
\]

\subsubsection{Почему дельта действует как
подстановка}\label{ux43fux43eux447ux435ux43cux443-ux434ux435ux43bux44cux442ux430-ux434ux435ux439ux441ux442ux432ux443ux435ux442-ux43aux430ux43a-ux43fux43eux434ux441ux442ux430ux43dux43eux432ux43aux430}

Для простой функции это сразу видно:

\[
\int \sum_j c_j\mathbf 1_{A_j}(x)\,\delta_a(dx)
=
\sum_j c_j\mathbf 1_{A_j}(a).
\]

Правая часть равна значению простой функции в точке \(a\). Переходом к
пределу получаем для ограниченной измеримой функции:

\[
\int f\,d\delta_a=f(a).
\]

\subsubsection{Почему характеристическая функция -
Фурье-образ}\label{ux43fux43eux447ux435ux43cux443-ux445ux430ux440ux430ux43aux442ux435ux440ux438ux441ux442ux438ux447ux435ux441ux43aux430ux44f-ux444ux443ux43dux43aux446ux438ux44f---ux444ux443ux440ux44cux435-ux43eux431ux440ux430ux437}

Если \(X\sim\mu\), то по формуле замены переменной:

\[
\mathbb E e^{i\langle u,X\rangle}
=
\int e^{i\langle u,x\rangle}\,\mu(dx).
\]

Левая часть - характеристическая функция, правая - преобразование Фурье
меры.

\subsubsection{Почему при независимости характеристики
перемножаются}\label{ux43fux43eux447ux435ux43cux443-ux43fux440ux438-ux43dux435ux437ux430ux432ux438ux441ux438ux43cux43eux441ux442ux438-ux445ux430ux440ux430ux43aux442ux435ux440ux438ux441ux442ux438ux43aux438-ux43fux435ux440ux435ux43cux43dux43eux436ux430ux44eux442ux441ux44f}

Если \(X\) и \(Y\) независимы, то

\[
\varphi_{X+Y}(u)
=
\mathbb E e^{i\langle u,X+Y\rangle}
=
\mathbb E\left(e^{i\langle u,X\rangle}e^{i\langle u,Y\rangle}\right).
\]

По независимости:

\[
\varphi_{X+Y}(u)
=
\mathbb E e^{i\langle u,X\rangle}
\mathbb E e^{i\langle u,Y\rangle}
=
\varphi_X(u)\varphi_Y(u).
\]

\subsection{7. Что написать на
доске}\label{ux447ux442ux43e-ux43dux430ux43fux438ux441ux430ux442ux44c-ux43dux430-ux434ux43eux441ux43aux435}

\begin{enumerate}
\def\labelenumi{\arabic{enumi}.}
\tightlist
\item
  Мера как функционал:
\end{enumerate}

\[
\langle\mu,f\rangle=\int f(x)\,\mu(dx).
\]

\begin{enumerate}
\def\labelenumi{\arabic{enumi}.}
\setcounter{enumi}{1}
\tightlist
\item
  Конечномерное распределение:
\end{enumerate}

\[
\langle P_{t_1,\ldots,t_n},f\rangle
=
\mathbb E f(X_{t_1},\ldots,X_{t_n}).
\]

\begin{enumerate}
\def\labelenumi{\arabic{enumi}.}
\setcounter{enumi}{2}
\tightlist
\item
  Дельта Дирака:
\end{enumerate}

\[
\delta_a(B)=\mathbf 1_B(a),
\qquad
\langle\delta_a,f\rangle=f(a).
\]

\begin{enumerate}
\def\labelenumi{\arabic{enumi}.}
\setcounter{enumi}{3}
\tightlist
\item
  Дискретный закон:
\end{enumerate}

\[
\mu=\sum_kp_k\delta_{x_k}.
\]

\begin{enumerate}
\def\labelenumi{\arabic{enumi}.}
\setcounter{enumi}{4}
\tightlist
\item
  Фурье-образ меры:
\end{enumerate}

\[
\widehat\mu(u)=\int e^{i\langle u,x\rangle}\,\mu(dx).
\]

\begin{enumerate}
\def\labelenumi{\arabic{enumi}.}
\setcounter{enumi}{5}
\tightlist
\item
  Характеристическая функция:
\end{enumerate}

\[
\varphi_X(u)=\widehat{P_X}(u).
\]

\begin{enumerate}
\def\labelenumi{\arabic{enumi}.}
\setcounter{enumi}{6}
\tightlist
\item
  Слабая сходимость:
\end{enumerate}

\[
\mu_n\Rightarrow\mu
\Longleftrightarrow
\int f\,d\mu_n\to\int f\,d\mu
\quad(f\in C_b).
\]

\subsection{8. Типичные дополнительные
вопросы}\label{ux442ux438ux43fux438ux447ux43dux44bux435-ux434ux43eux43fux43eux43bux43dux438ux442ux435ux43bux44cux43dux44bux435-ux432ux43eux43fux440ux43eux441ux44b}

\textbf{Что значит рассматривать распределение как обобщенную функцию?}

Это значит рассматривать меру \(\mu\) по ее действию на тестовые
функции:

\[
f\mapsto\langle\mu,f\rangle=\int f\,d\mu.
\]

Даже если у распределения нет плотности, такой функционал определен.

\textbf{Почему берут ограниченные непрерывные тестовые функции?}

Ограниченность гарантирует конечность интеграла по вероятностной мере.
Непрерывность нужна для слабой сходимости: распределения считаются
близкими, если интегралы от всех \(f\in C_b\) близки.

\textbf{Чем дельта Дирака отличается от обычной функции?}

Строго \(\delta_a\) - это мера или функционал, а не обычная функция. Ее
смысл:

\[
\int f\,d\delta_a=f(a).
\]

Нельзя обращаться с ней как с обычной плотностью относительно меры
Лебега.

\textbf{Как записать дискретное распределение через дельта-меры?}

Если

\[
P\{X=x_k\}=p_k,
\]

то

\[
P_X=\sum_kp_k\delta_{x_k}.
\]

\textbf{Что такое преобразование Фурье меры?}

Это функция

\[
\widehat\mu(u)=\int e^{i\langle u,x\rangle}\,\mu(dx).
\]

Для вероятностной меры она всегда определена, потому что модуль
экспоненты равен \(1\).

\textbf{Почему характеристическая функция - это Фурье-образ
распределения?}

Если \(X\sim\mu\), то

\[
\varphi_X(u)
=
\mathbb E e^{i\langle u,X\rangle}
=
\int e^{i\langle u,x\rangle}\,\mu(dx)
=
\widehat\mu(u).
\]

\textbf{Что такое слабая сходимость в этом языке?}

Это сходимость мер как функционалов:

\[
\mu_n\Rightarrow\mu
\]

если для всех \(f\in C_b\):

\[
\langle\mu_n,f\rangle\to\langle\mu,f\rangle.
\]

\textbf{Какой Фурье-образ у дельта-меры?}

\[
\widehat{\delta_a}(u)=e^{i\langle u,a\rangle}.
\]

\textbf{Как связаны свертка и преобразование Фурье?}

Свертка мер переходит в произведение Фурье-образов:

\[
\widehat{\mu*\nu}(u)=\widehat\mu(u)\widehat\nu(u).
\]

Вероятностно это означает, что характеристическая функция суммы
независимых величин равна произведению характеристических функций.

\textbf{Какая главная ошибка в этом билете?}

Путать дельта-функцию с обычной функцией и пытаться использовать
плотность там, где распределение атомарное.

\subsection{9. Типичные
ошибки}\label{ux442ux438ux43fux438ux447ux43dux44bux435-ux43eux448ux438ux431ux43aux438}

\begin{itemize}
\tightlist
\item
  Путать дельта-меру с обычной функцией.
\item
  Писать плотность для атомарного распределения без оговорки о
  дельта-мерах.
\item
  Забывать, что \(\delta_a\) действует на функцию как подстановка.
\item
  Путать функцию распределения \(F(x)\) и распределение как меру
  \(\mu\).
\item
  Путать преобразование Фурье плотности и преобразование Фурье меры:
  второе работает и без плотности.
\item
  Забывать знак в соглашении Фурье; в вероятности обычно используют
  \(e^{i\langle u,x\rangle}\).
\item
  Считать, что слабая сходимость проверяется по индикаторам всех
  множеств; стандартная формулировка здесь через \(C_b\).
\end{itemize}

\subsection{10. Связи с другими
билетами}\label{ux441ux432ux44fux437ux438-ux441-ux434ux440ux443ux433ux438ux43cux438-ux431ux438ux43bux435ux442ux430ux43cux438}

\begin{itemize}
\tightlist
\item
  Билет 01: мера Дирака как пример меры.
\item
  Билет 06: распределения случайных величин и конечномерные
  распределения.
\item
  Билет 08: характеристические функции как преобразования Фурье
  распределений.
\item
  Билет 09: независимость и произведение характеристических функций.
\item
  Билет 13: моменты и ковариации через производные характеристической
  функции.
\item
  Билет 15: согласованные конечномерные распределения задают процесс.
\item
  Билет 18: гауссовские распределения удобно задаются
  характеристическими функциями.
\item
  Билет 19: для сумм независимых шагов характеристические функции
  перемножаются.
\end{itemize}

\subsection{11. Минимум на
удовлетворительно}\label{ux43cux438ux43dux438ux43cux443ux43c-ux43dux430-ux443ux434ux43eux432ux43bux435ux442ux432ux43eux440ux438ux442ux435ux43bux44cux43dux43e}

Нужно сказать:

\begin{enumerate}
\def\labelenumi{\arabic{enumi}.}
\tightlist
\item
  Распределение \(\mu\) можно рассматривать как функционал:
\end{enumerate}

\[
\langle\mu,f\rangle=\int f\,d\mu.
\]

\begin{enumerate}
\def\labelenumi{\arabic{enumi}.}
\setcounter{enumi}{1}
\tightlist
\item
  Для конечномерного распределения процесса:
\end{enumerate}

\[
\langle P_{t_1,\ldots,t_n},f\rangle
=
\mathbb E f(X_{t_1},\ldots,X_{t_n}).
\]

\begin{enumerate}
\def\labelenumi{\arabic{enumi}.}
\setcounter{enumi}{2}
\tightlist
\item
  Дельта Дирака:
\end{enumerate}

\[
\langle\delta_a,f\rangle=f(a).
\]

\begin{enumerate}
\def\labelenumi{\arabic{enumi}.}
\setcounter{enumi}{3}
\tightlist
\item
  Преобразование Фурье меры:
\end{enumerate}

\[
\widehat\mu(u)=\int e^{i\langle u,x\rangle}\,\mu(dx).
\]

\begin{enumerate}
\def\labelenumi{\arabic{enumi}.}
\setcounter{enumi}{4}
\tightlist
\item
  Для распределения случайного вектора это характеристическая функция.
\end{enumerate}

\subsection{12. Минимум на
хорошо}\label{ux43cux438ux43dux438ux43cux443ux43c-ux43dux430-ux445ux43eux440ux43eux448ux43e}

Помимо минимума, нужно объяснить:

\begin{itemize}
\tightlist
\item
  почему ограниченность тестовых функций гарантирует конечность
  интеграла;
\item
  почему дельта Дирака не является обычной плотностью;
\item
  как записать дискретное распределение через сумму дельта-мер;
\item
  как записать абсолютно непрерывное распределение через плотность;
\item
  что слабая сходимость означает сходимость интегралов по \(C_b\);
\item
  почему характеристическая функция однозначно задает распределение.
\end{itemize}

\subsection{13. Что добавить для
отлично}\label{ux447ux442ux43e-ux434ux43eux431ux430ux432ux438ux442ux44c-ux434ux43bux44f-ux43eux442ux43bux438ux447ux43dux43e}

Для сильного ответа стоит добавить:

\begin{itemize}
\tightlist
\item
  связь с конечномерными распределениями процесса;
\item
  положительность и нормированность функционала
  \(\mu:f\mapsto\int f\,d\mu\);
\item
  примеры: константа, Бернулли, плотность;
\item
  преобразование Фурье дельта-меры;
\item
  свертку и произведение характеристических функций;
\item
  предупреждение про соглашение знака в преобразовании Фурье;
\item
  связь со слабой сходимостью и ЦПТ.
\end{itemize}

\subsection{14. Устная версия ответа на 5
минут}\label{ux443ux441ux442ux43dux430ux44f-ux432ux435ux440ux441ux438ux44f-ux43eux442ux432ux435ux442ux430-ux43dux430-5-ux43cux438ux43dux443ux442}

\begin{enumerate}
\def\labelenumi{\arabic{enumi}.}
\tightlist
\item
  Конечномерное распределение процесса - это мера на \(\mathbb R^n\).
\item
  Эту меру можно рассматривать как функционал на тестовых функциях:
\end{enumerate}

\[
\langle\mu,f\rangle=\int f\,d\mu.
\]

\begin{enumerate}
\def\labelenumi{\arabic{enumi}.}
\setcounter{enumi}{2}
\tightlist
\item
  Для процесса:
\end{enumerate}

\[
\langle P_{t_1,\ldots,t_n},f\rangle
=
\mathbb E f(X_{t_1},\ldots,X_{t_n}).
\]

\begin{enumerate}
\def\labelenumi{\arabic{enumi}.}
\setcounter{enumi}{3}
\tightlist
\item
  Такой функционал линейный, положительный и нормированный.
\item
  Дельта Дирака в точке \(a\):
\end{enumerate}

\[
\delta_a(B)=\mathbf 1_B(a),
\qquad
\langle\delta_a,f\rangle=f(a).
\]

\begin{enumerate}
\def\labelenumi{\arabic{enumi}.}
\setcounter{enumi}{5}
\tightlist
\item
  Дискретное распределение записывается как сумма дельт:
\end{enumerate}

\[
\mu=\sum_kp_k\delta_{x_k}.
\]

\begin{enumerate}
\def\labelenumi{\arabic{enumi}.}
\setcounter{enumi}{6}
\tightlist
\item
  Если есть плотность \(p\), то
\end{enumerate}

\[
\langle\mu,f\rangle=\int f(x)p(x)\,dx.
\]

\begin{enumerate}
\def\labelenumi{\arabic{enumi}.}
\setcounter{enumi}{7}
\tightlist
\item
  Преобразование Фурье меры:
\end{enumerate}

\[
\widehat\mu(u)=\int e^{i\langle u,x\rangle}\,\mu(dx).
\]

\begin{enumerate}
\def\labelenumi{\arabic{enumi}.}
\setcounter{enumi}{8}
\tightlist
\item
  Для вероятностной меры это характеристическая функция:
\end{enumerate}

\[
\varphi_X(u)=\widehat{P_X}(u).
\]

\begin{enumerate}
\def\labelenumi{\arabic{enumi}.}
\setcounter{enumi}{9}
\tightlist
\item
  У дельта-меры:
\end{enumerate}

\[
\widehat{\delta_a}(u)=e^{i\langle u,a\rangle}.
\]

\begin{enumerate}
\def\labelenumi{\arabic{enumi}.}
\setcounter{enumi}{10}
\tightlist
\item
  Слабая сходимость означает сходимость интегралов от всех ограниченных
  непрерывных тестовых функций.
\item
  Главная ошибка - считать \(\delta\) обычной функцией и пытаться всюду
  писать плотности.
\end{enumerate}

\subsection{15. Если осталось 2 минуты перед
экзаменом}\label{ux435ux441ux43bux438-ux43eux441ux442ux430ux43bux43eux441ux44c-2-ux43cux438ux43dux443ux442ux44b-ux43fux435ux440ux435ux434-ux44dux43aux437ux430ux43cux435ux43dux43eux43c}

Запомнить три формулы.

Мера как функционал:

\[
\langle\mu,f\rangle=\int f\,d\mu.
\]

Дельта Дирака:

\[
\langle\delta_a,f\rangle=f(a).
\]

Фурье-образ меры:

\[
\widehat\mu(u)=\int e^{i\langle u,x\rangle}\,\mu(dx).
\]

Для распределения \(X\):

\[
\widehat{P_X}(u)=\varphi_X(u).
\]

Главная ловушка: \(\delta_a\) - мера/функционал, а не обычная функция
плотности.

\newpage

\end{document}
